\begin{abstract}
This iteration combines recent ProveIt constructions with a later-stage application of
paired sign conditions. If two integers have a known nonnegative sum,
one nonnegativity certificate for their product enforces both signs.
Applied after a carefully ordered power-recovery argument, this gives a
direct quartic in 489 unrestricted integers, compared with the incoming
517-variable predecessor. The variable saving is accompanied by larger
coefficient height, not an improvement in every size measure.

A second centering gives a simpler compact endpoint: fifteen unrestricted
integers, exact degree 48 and 108 arithmetic operations, compared with
degree 56 and 110 operations. Coprime grouping supplies new eleven-integer
tradeoffs, while auxiliary Pell growth removes a transport majorant from
ten-positive-variable equations. The smallest operation count in that
ten-variable family is 89 at degree 288; a degree-254 alternative uses
97 operations. All conclusions concern ordinary mathematical reductions
with explicit inherited logical calibration, not new end-to-end kernel
certificates or a global optimum.
\end{abstract}

\section*{Principal results}
\begin{center}
\begin{tabular}{@{}lrrl@{}}
\toprule
Domain and variables & Degree & Operations & Change or tradeoff\\
\midrule
$\mathbb Z^{15}$ &48&108&Previously $56/110$\\
$\mathbb Z^{12}$ &500&139&Same polynomial, previously 148 gates\\
$\mathbb Z^{11}$ &688&124&Same polynomial, previously 126 gates\\
$\mathbb Z^{11}$ &820&119&New intermediate tradeoff\\
$\mathbb Z^{11}$ &964&118&New intermediate tradeoff\\
$\mathbb Z^{11}$ &1632&112&Same polynomial, previously 114 gates\\
$\mathbb Z_{>0}^{10}$ &288&89&Shorter circuit, larger degree\\
$\mathbb Z_{>0}^{10}$ &262&94&Previously $266/96$\\
$\mathbb Z_{>0}^{10}$ &254&97&Incoming predecessor $254/99$\\
$\mathbb Z^{489}$ &4&---&Direct source, previously 517 variables\\
\bottomrule
\end{tabular}
\end{center}
Here $d/c$ means degree $d$ and $c$ live binary operations. The dash does
not mean zero operations: the direct quartic is compared by its explicit
quadratic residuals, expanded terms and coefficient height instead.
Additional tradeoffs and exact files are given in the relevant sections.
The existing nine-variable positive constructions and shorter circuits
with more variables remain separate comparison points. No claim is made
that 89 operations or ten variables is a universal minimum.

The direct comparison is particularly important:
\begin{center}
\begin{tabular}{@{}rrrrr@{}}
\toprule
Integer variables & Degree & Residuals & Terms & Height\\
\midrule
517 (predecessor)&4&310&6774&3904267\\
493 (paired signs)&4&310&7174&10945213\\
489 (square anchor)&4&308&6905&11302118\\
\bottomrule
\end{tabular}
\end{center}
The twelve post-power sign pairs and square anchor extend a contemporaneous
upstream six-pair construction; the principle itself is credited, not
claimed as new. The fifteen-variable result is an explicit zero-set
bijection. Other entries combine existing reductions or optimize evaluators.
